We begin with a rather complicated identity for the joint distribution
of the quantities $C_{I_i}$.

\begin{lem}\label{lem:CIjmj}
Let $I_1,\ldots,I_r$ be disjoint subsets of $[n]$ and $m_1,\ldots,m_r$
be non-negative integers.  Denote $I_0 = [n] \setminus (I_1\cup \cdots \cup I_r)$.
Let $\cT$ be the set of indices $i$ with $m_i>0$, together
with the number 0 if $I_0$ is nonempty.  Then
\[
  \PR\big( C_{I_j}(\sigma)=m_j\; (1\le j\le r) \big)
  = \frac{1}{n}\sum_{t\in \cT} \sum_{h \in I_t} 
  \ssum{b_1,\dots, b_n \geq 0 \\ b_1 + 2b_2 + \cdots + nb_n = n - h \\ \sum_{i\in I_j} b_i =m_j-\one(t=j),\  (1\le j\le r) } 
  \prod_{i=1}^n \frac{(1/i)^{b_i}}{b_i!}.
\]
\end{lem}

\begin{proof}
 Evidently
\[
  n \# \{\sigma \in \cS_n:  C_{I_1}(\sigma)=m_1,\ldots,C_{I_r}(\sigma)=m_r\} = \ssum{\sigma \in \cS_n \\ C_{I_j}(\sigma)=m_j \; (1\le j\le r) } \;\; \ssum{\alpha|\sigma \\ \a \text{ a cycle}} |\a|.
\]
Write $\sigma=\a \beta$ and let $h=|\a|$.   
 Thus, for some $t\in \cT$,  we have $|\a|=h\in I_t$
 and
\[
(C_{I_1}(\beta),\ldots,C_{I_r}(\beta))=(m_1-\one(t=1),\ldots,m_r-\one(t=r)).
\]
It is permissible to think of $\beta\in \cS_{n-h}$ and thus
\begin{align*}
  n  \# \{\sigma \in \cS_n: C_{I_1}(\sigma)=m_1,\ldots,C_{I_r}(\sigma)=m_r\}
  &= \sum_{t\in \cT} \sum_{h\in I_t} \; \ssum{\a\in\cS_n, |\a|=h \\ \a\text{ a cycle}} h \ssum{\beta\in \cS_{n-h}\\ C_{I_i}(\beta)=m_i-\one(t=i), (1\le i\le r)} 1  \\
  &=  \sum_{t\in \cT} \sum_{h\in I_t} \frac{n!}{(n-h)!} \ssum{\beta\in \cS_{n-h} \\ C_{I_i}(\beta)=m_i-\one(t=i), (1\le i\le r) } 1.
\end{align*}
Now subdivide the sum according the cycle type $(b_1,\ldots,b_{n})$ of
the permutation $\beta$, use Cauchy's formula (Thm. \ref{thm:Cauchy}) to 
count such permutations for each type, and divide by $n$.  The desired identity follows.
\end{proof}




\begin{proof}[Proof of Theorem \ref{cycles_sets}]
The right side in Lemma \ref{lem:CIjmj} is at most
 \[
 \frac{1}{n} \sum_{t\in \cT}
  \ssum{b_1,\dots, b_n \geq 0 \\ \sum_{i\in I_j} b_i =m_j-\one(t=j) \; (1\le j\le r)  } \prod_i \frac{(1/i)^{b_i}}{b_i!} := \frac{Y}{n},
\]
  say.   By the multinomial theorem,
  \begin{align*}
 Y & =    \sum_{t\in \cT} \ssum{b_i\ge 0 \; (i\in I_1\cup \cdots \cup I_r) \\ \sum_{i\in I_j} b_i = m_j-\one(t=j) \; (1\le j\le r) } 
  \frac{1}{\prod_{i\in I_1\cup \cdots \cup I_r} b_i! i^{b_i}} \sum_{b_i\ge 0 \; (i\in I_0)} \frac1{\prod_{i\in I_0} b_i! i^{b_i}}\\
  &=  \sum_{t\in \cT} \frac{m_t}{H(I_t)}\prod_{j=1}^r \frac{H(I_j)^{m_j}}{m_j!} \er^{H(I_0)}.
\end{align*}
 The claimed bound now follows from
$H(I_0) = H_n - H(I_1)-\cdots-H(I_r)$.
\end{proof}
